\section{Reliability, Hallucination, and Legal Citation Verification}
\label{sec:reliability}

\subsection{From relevant retrieval to reliable legal answers}

Retrieval-augmented generation (RAG) has become an important approach to reducing the dependence of
large language models (LLMs) on parametric knowledge alone. By retrieving external documents at
inference time, a RAG system can provide answers grounded in legal statutes, regulations, judgments,
and other authoritative materials. However, access to external evidence does not automatically make an
answer reliable. A system may retrieve a relevant document but misinterpret its legal meaning, overlook
an exception, cite the wrong passage, or generate a conclusion that the retrieved evidence does not
support.

This distinction is particularly important in law because legal correctness depends on more than
textual similarity. An answer may accurately reproduce a sentence from a judgment while misrepresenting
its holding. It may identify the correct statute but overlook an amendment, jurisdictional restriction,
procedural condition, or subsequent decision that affects its application. Consequently, reliability
must be evaluated across the complete pipeline, from evidence retrieval to legal interpretation and
citation generation.

The reliability problem can be understood through three related questions:
\begin{enumerate}[leftmargin=1.4em, itemsep=2pt, topsep=2pt]
\item \textbf{Was the right evidence retrieved?} The system must locate the relevant legal provision,
passage, or judgment.
\item \textbf{Does the evidence support the claim?} The generated statement must follow from the cited
material rather than merely discuss a similar topic.
\item \textbf{Is the claim legally applicable?} The system must account for relevant jurisdiction,
temporal validity, procedural context, and legal authority.
\end{enumerate}

These questions are related but not interchangeable. A system can perform well on retrieval metrics and
still generate unsupported conclusions. Similarly, a response containing accurate quotations may remain
misleading if it applies them to the wrong legal issue.

\subsection{Types of reliability failures}

Legal AI reliability is not a single measurable property. It includes several failure modes that
require different detection and mitigation strategies, as summarized in Table~\ref{tab:failures}.

\begin{table*}[tp]
\centering
\caption{Reliability failure modes in legal AI and possible mitigations.}
\label{tab:failures}
\footnotesize
\setlength{\tabcolsep}{5pt}
\renewcommand{\arraystretch}{1.3}
\begin{tabularx}{\textwidth}{>{\raggedright\arraybackslash}p{2.5cm} >{\raggedright\arraybackslash}X >{\raggedright\arraybackslash}X >{\raggedright\arraybackslash}X}
\toprule
\textbf{Failure type} & \textbf{Description} & \textbf{Illustrative example} & \textbf{Possible mitigation} \\
\midrule
Fabricated authority & The system invents a case, statute, provision, or citation. & Referring to a nonexistent judgment. & Citation resolution against authoritative databases. \\
Incorrect citation & A real source is cited, but the cited passage does not support the associated claim. & Citing a procedural provision as authority for a substantive conclusion. & Claim-to-passage entailment checking. \\
Unsupported claim & The response makes a factual or legal assertion that is not established by retrieved evidence. & Stating that a remedy is available without retrieving the relevant conditions. & Evidence coverage checks and abstention. \\
Misinterpretation & The source is real, but its legal meaning is misrepresented. & Treating an observation as the court's binding holding. & Role-aware judgment analysis and contextual verification. \\
Incomplete evidence & A relevant exception, counter-authority, or condition is omitted. & Explaining a general rule without its statutory exception. & Query expansion and retrieval of related provisions. \\
Outdated law & The system relies on a superseded or inapplicable version of a legal source. & Applying an earlier statutory amendment. & Effective-date metadata and version-aware retrieval. \\
Jurisdictional mismatch & The system applies law from a different jurisdiction. & Using an Indian legal rule to answer a Pakistani-law question. & Jurisdiction filtering and authority classification. \\
Contradictory evidence & The system fails to reconcile conflicting authorities. & Presenting one judgment without acknowledging a relevant later ruling. & Authority-aware ranking and explicit conflict reporting. \\
\bottomrule
\end{tabularx}
\end{table*}

The categories overlap. For example, a jurisdictional mismatch may also produce an unsupported claim,
while outdated law can lead to incorrect citations. Nevertheless, distinguishing the underlying failure
is useful because a single safeguard cannot address every problem.

\subsection{Evidence grounding and citation correctness}

A citation should not be treated as reliable merely because it points to an existing document. Legal
citation verification requires checking both the existence of the source and its relationship to the
claim. A practical verification process separates four operations:
\begin{itemize}[leftmargin=1.4em, itemsep=2pt, topsep=2pt]
\item \textbf{Citation resolution:} Determine whether the cited statute, case, section, paragraph, or
other authority exists and can be located.
\item \textbf{Passage verification:} Confirm that the cited passage contains the relevant legal
language or reasoning.
\item \textbf{Claim-support verification:} Assess whether the passage actually supports the claim,
rather than simply sharing vocabulary with it.
\item \textbf{Applicability verification:} Check whether the authority is relevant to the jurisdiction,
date, issue, and procedural context of the question.
\end{itemize}

This process suggests that citation verification should operate at the level of individual claims, not
merely at the level of the complete response. A long answer may contain several supported statements and
one critical unsupported conclusion. Evaluating the answer as a single unit can obscure that distinction.

LegalBench-RAG addresses an important part of this problem by focusing on retrieving precise legal text
segments rather than relying solely on document-level retrieval. Its benchmark design reinforces the
importance of locating evidence at a sufficiently granular level for downstream legal tasks
\citep{pipitone2024}. However, precise retrieval and legal entailment remain separate requirements:
finding the correct passage does not itself establish that the generated interpretation is valid.

\citet{dahl2024}, in their discussion of fabricated legal authorities and the broader reliability of
generative AI in law, and \citet{magesh2025}, in their assessment of leading AI legal research tools,
further motivate evaluating the output rather than assuming that fluent language or a citation list
indicates correctness. These studies should not be interpreted as establishing one universal
hallucination rate for all legal AI systems. Results depend on the tools, tasks, jurisdictions, test
sets, and definitions of error used in each evaluation.

\subsection{Evaluating reliability: beyond answer accuracy}

A robust evaluation framework should distinguish retrieval quality, citation quality, factual or legal
support, and operational performance. No single metric captures all of these dimensions, as
Table~\ref{tab:releval} shows.

\begin{table*}[tp]
\centering
\caption{Dimensions for evaluating legal AI reliability.}
\label{tab:releval}
\footnotesize
\setlength{\tabcolsep}{5pt}
\renewcommand{\arraystretch}{1.3}
\begin{tabularx}{\textwidth}{>{\raggedright\arraybackslash}p{2.6cm} >{\raggedright\arraybackslash}p{2.9cm} >{\raggedright\arraybackslash}X >{\raggedright\arraybackslash}X}
\toprule
\textbf{Evaluation dimension} & \textbf{Example metric} & \textbf{What it measures} & \textbf{Important limitation} \\
\midrule
Retrieval effectiveness & Recall@$k$ & Whether relevant evidence appears among the top $k$ retrieved results. & Does not measure whether the final answer uses that evidence correctly. \\
Ranking quality & MRR or nDCG & How effectively the system ranks relevant passages. & Depends on relevance judgments and benchmark design. \\
Citation validity & Citation validity rate & Proportion of generated citations that resolve to real sources. & A real citation may still be irrelevant or misleading. \\
Citation support & Claim--citation support rate & Proportion of claims adequately supported by their cited evidence. & Requires careful annotation or reliable support verification. \\
Evidence coverage & Evidence coverage rate & Extent to which material claims have supporting evidence. & Does not establish that the evidence is legally correct or complete. \\
Unsupported content & Unsupported-claim rate & Proportion of claims lacking adequate evidential support. & Results vary with claim segmentation and annotation policy. \\
Legal correctness & Expert-assessed correctness & Whether the answer's legal analysis is judged correct under the task conditions. & Requires qualified evaluators and explicit jurisdictional criteria. \\
Calibration & Confidence--accuracy alignment & Whether confidence corresponds to observed correctness. & Model-reported confidence can be poorly calibrated. \\
Efficiency & Latency and cost per query & Operational expense of producing an answer. & Low cost or latency does not imply legal reliability. \\
\bottomrule
\end{tabularx}
\end{table*}

Metric definitions should be fixed before evaluation. For example, citation validity rate and citation
support rate should not be used interchangeably. A citation can be valid in the first sense and invalid
in the second. Likewise, an unsupported-claim rate should specify whether the denominator consists of
all generated claims, all material claims, or all claims requiring citations.

Evaluation should also distinguish automatic checks from expert assessment. Automated tools can identify
broken links, malformed citations, missing sources, and some forms of semantic inconsistency. Human
legal reviewers remain important for assessing nuanced statutory interpretation, the relative authority
of cases, exceptions, and whether a response appropriately qualifies uncertainty. Ideally, evaluations
should report the annotation protocol, reviewer expertise, agreement between reviewers, and adjudication
procedure.

\subsection{Safeguards for legal AI applications}

Reliability is best addressed through multiple safeguards rather than a single verification component:
\begin{enumerate}[leftmargin=1.4em, itemsep=2pt, topsep=2pt]
\item \textbf{Authoritative source management.} Record provenance, jurisdiction, document identity,
effective dates, and amendments.
\item \textbf{Evidence retrieval and ranking.} Retrieve relevant passages using lexical, dense, hybrid,
or structure-aware methods.
\item \textbf{Claim-level generation.} Associate material legal claims with specific supporting
passages.
\item \textbf{Verification and conflict checks.} Resolve citations, test claim support, and identify
missing or conflicting authorities.
\item \textbf{Qualified response or abstention.} Answer with traceable evidence, state uncertainty, or
decline to make an unsupported conclusion.
\end{enumerate}
\textit{This is a conceptual reliability workflow synthesized for this survey; it is not a claim that
every reviewed system implements all five stages.}

The workflow highlights several design principles. First, provenance and temporal metadata should be
preserved during ingestion rather than added after generation. Second, verification should be applied
to material claims and their supporting passages. Third, the system should be allowed to abstain when
evidence is missing or contradictory. Finally, high-impact legal applications should preserve access to
the underlying sources and support review by a qualified human.

Agentic architectures can coordinate these safeguards through retrieval, verification, and revision
steps. However, adding more agents does not automatically improve correctness. If the agents rely on the
same incomplete corpus or repeat the same unsupported assumption, orchestration may reproduce rather
than correct the error. Reliability therefore depends on independent evidence checks, explicit
verification criteria, and measurable performance improvements.

\subsection{Section synthesis}

The central lesson is that legal AI reliability must be evaluated as an evidence-to-claim problem.
Retrieval quality is necessary, but it is not sufficient. A dependable system must retrieve appropriate
authority, preserve its context, generate claims that follow from the evidence, verify citations, and
recognize when legal applicability remains uncertain. These requirements motivate the comparative
analysis in the next section, where contemporary systems are assessed according to their evidence
sources, retrieval design, legal reasoning, verification mechanisms, and evaluation practices.

\subsection{Implications for evaluation and reproducibility}

The literature also highlights a methodological problem: reliability results cannot be compared fairly
unless the underlying evaluation conditions are reported. A hallucination rate obtained from open-ended
questions about legal doctrine is not directly comparable with one obtained from citation verification,
statutory question answering, or retrieval of annotated passages.

For example, \citet{magesh2025} evaluated commercial legal research tools and found that hallucinations
persisted despite the use of RAG. Their findings challenge claims that retrieval alone can guarantee
hallucination-free outputs. \citet{dahl2024} likewise documented substantial hallucination problems in
general-purpose LLMs answering legal questions. A reproducible legal AI evaluation should therefore
disclose:
\begin{itemize}[leftmargin=1.4em, itemsep=2pt, topsep=2pt]
\item the jurisdiction, legal domain, and temporal scope of the corpus;
\item the source of questions and the criteria used to construct the test set;
\item the operational definition of hallucination and unsupported claims;
\item the annotation instructions, reviewer expertise, and agreement statistics;
\item the retrieval configuration, model version, prompts, and verification settings;
\item the baseline systems, uncertainty estimates, and statistical tests where appropriate.
\end{itemize}

These requirements are especially important when evaluating improvements from hybrid retrieval,
reranking, knowledge graphs, or agentic verification. Each component should be assessed through
controlled comparisons or ablation studies, rather than attributing all observed improvements to the
complete architecture.

The resulting research priority is clear: legal AI needs not only more capable models, but also
transparent evidence pipelines and evaluations that measure whether answers are actually supported by
applicable legal authority.
